% Part: methods
% Chapter: induction
% Section: strong-induction

\documentclass[../../../include/open-logic-section]{subfiles}

\begin{document}

\olfileid{mth}{ind}{str}

\olsection{强归纳法}

在前面讨论的归纳法原理中，我们证明了$P(0)$，并证明若$P(k)$则$P(k+1)$。在第二部分，我们假设$P(k)$为真，并利用该假设证明$P(k+1)$。当然，等价地，我们也可以假设$P(k-1)$并用它证明$P(k)$——关键在于，我们能够从任意数到其后继进行推理；即在假设命题对其前驱成立的情况下，能够证明该命题对该数本身成立。

归纳法原理有一个变体：我们不仅假设命题对$k$的前驱$k-1$成立，而且假设它对所有小于$k$的数成立，并利用这一假设来确立命题对$k$成立。这同样能得出对所有$n \in \Nat$都有$P(n)$。因为一旦确立了$P(0)$，我们也就确立了$P$对所有小于$1$的数都成立。如果我们知道“若对所有$l<k$都有$P(l)$，则$P(k)$”，那么特别地，当$k=1$时这也成立。于是我们可以得出$P(1)$。由此我们证明了$P(0)$和$P(1)$，亦即对所有$l<2$都有$P(l)$；既然我们还有条件句“若对所有$l<2$都有$P(l)$，则$P(2)$”，我们就可以得出$P(2)$，依此类推。

事实上，如果我们能确立一般条件句“对所有$k$，若对所有$l<k$都有$P(l)$，则$P(k)$”，我们就不再需要单独确立$P(0)$了，因为它可以从中得出。因为要知道，像“对所有$l<k$，$P(l)$”这样的全称命题，当不存在任何$l<k$时，它也是真的。这是空量化的一个例子：如果不存在$A$，则“所有$A$都是$B$”为真；如果没有$x$满足$!A(x)$，则$\lforall[x][(!A(x) \lif !B(x))]$为真。在此情形下，形式化的版本应为“$\lforall[l][(l < k \lif P(l))]$”——而当不存在任何$l < k$时，该命题即为真。而当$k=0$时恰好就是这种情形：没有$l<0$，因此“对所有$l<0$，$P(0)$”为真，无论$P$是什么。于是，对“若对所有$l<k$都有$P(l)$，则$P(k)$”的证明便自动确立了$P(0)$。

当确立命题对$k$成立不能仅仅依赖于它对$k-1$成立，而可能需要假设它对一个或多个$l<k$为真时，这种变体就很有用。

\end{document}